% from ega4-2-fr.tex
\oldpage[IV-2]{41}
fermées de $X$). Dire que $\mathcal O_X$ est réduit en un point $x$ équivaut à
dire que $X$ est \emph{réduit au point $x$}.

\begin{proposition}[3.2.3]
\label{IV.2.3.2.3-fr}
Soient $X$ un préschéma localement noethérien, $U$ un ouvert dans $X$,
$\mathcal F$ un $\mathcal O_X$-Module cohérent tel que
$U\cap\operatorname{Supp}(\mathcal F)$ soit dense dans
$\operatorname{Supp}(\mathcal F)$. Les conditions suivantes sont
équivalentes~:
\begin{enumerate}
\item[\emph{a)}] $\mathcal F$ est réduit.
\item[\emph{b)}] $\mathcal F|U$ est réduit, et $\mathcal F$ est sans cycle
premier associé immergé.
\item[\emph{c)}] Il existe un sous-préschéma fermé réduit $X'$ de $X$ et un
$\mathcal O_{X'}$-Module cohérent $\mathcal F'$ sans torsion et de rang $1$
sur toute composante irréductible de $X'$, tels que si $j:X'\to X$ est
l'injection canonique, on ait $j_*(\mathcal F')=\mathcal F$.
\end{enumerate}
En outre, lorsqu'il en est ainsi le sous-préschéma $X'$ est défini par l'Idéal
$\mathcal J$ de $\mathcal O_X$ annulateur de $\mathcal F$.
\end{proposition}

Pour voir que \emph{c)} implique \emph{a)}, on peut, en vertu de (3.1.3), se
limiter au cas où $X'=X$ est intègre, de point générique $x$ et l'on peut en
outre supposer $X=\operatorname{Spec}(A)$ affine et
$\mathcal F=\widetilde M$, où $A$ est donc intègre et $M$ est un $A$-module
sans torsion de rang $1$~; l'annulateur de tout élément de $M$ étant alors
réduit à $0$, on a $\operatorname{Ass}(\mathcal F)=\{x\}$ et
$\mathcal F_x$ est isomorphe à $\mathcal O_x$, corps des fractions de $A$,
donc les conditions de la définition (3.2.2) sont satisfaites. Comme
l'existence de cycles premiers associés immergés est locale, il est clair que
si $\mathcal F$ n'a pas de tels cycles et si
$U\cap\operatorname{Supp}(\mathcal F)$ est dense dans
$\operatorname{Supp}(\mathcal F)$, on a
$\operatorname{Ass}(\mathcal F)=\operatorname{Ass}(\mathcal F|U)$, donc
\emph{a)} et \emph{b)} sont équivalentes. Si \emph{a)} est satisfaite,
prenons pour $X'$ le sous-préschéma fermé réduit de $X$ dont
$\operatorname{Supp}(\mathcal F)$ est l'espace sous-jacent (I, 5.2.1), et soit
$\mathcal F'=j^*(\mathcal F)$~; un point $x$ de
$\operatorname{Ass}(\mathcal F)$ est nécessairement point maximal de $X'$ et
comme $\operatorname{long}(\mathcal F_x)=1$, $\mathcal F'_x$ est isomorphe à
$\mathcal F_x$, donc au corps $k(x)=\mathcal O_{X',x}$, ce qui prouve que
$\mathcal F'$ est sans torsion et de rang $1$ (I, 7.4.6 et 7.4.2). Enfin, la
dernière assertion est triviale, puisque pour tout $y\in X'$, l'annulateur du
$\mathcal O_{X',y}$-Module $\mathcal F'_y$ est nul.

\begin{definition}[3.2.4]
\label{IV.2.3.2.4-fr}
Soient $X$ un préschéma localement noethérien, $\mathcal F$ un
$\mathcal O_X$-Module quasi-cohérent. On dit que $\mathcal F$ est
\emph{irredondant} si $\operatorname{Ass}(\mathcal F)$ est réduit à un seul
élément $x$~; si $\mathcal F$ est de type fini, on dit que $\mathcal F$ est
\emph{intègre} si de plus $\mathcal F$ est réduit (autrement dit si
$\operatorname{long}(\mathcal F_x)=1$). On dit qu'un sous-$\mathcal O_X$-Module
quasi-cohérent $\mathcal G$ d'un $\mathcal O_X$-Module quasi-cohérent
$\mathcal F$ est \emph{primaire dans $\mathcal F$} si
$\mathcal F/\mathcal G$ est irredondant.
\end{definition}

Pour un schéma affine $X$, cette définition donne en particulier la notion de
\emph{module intègre} sur un anneau noethérien $A$~; un $A$-module $M$ est dit
\emph{intègre} s'il est de type fini, si $M$ est \emph{primaire} (c'est-à-dire
que $\operatorname{Ass}(M)$ est réduit à un seul idéal premier
$\mathfrak p$) et si en outre
$\operatorname{long}_{A_\mathfrak p}(M_\mathfrak p)=1$. Revenant au cas d'un
préschéma localement noethérien quelconque $X$ et d'un
$\mathcal O_X$-Module cohérent $\mathcal F$, on dit que $\mathcal F$ est
\emph{intègre en un point} $x\in X$ si $\mathcal F_x$ est un
$\mathcal O_x$-module \emph{intègre}~: cela signifie encore que $x$
\emph{n'appartient qu'à un seul cycle premier associé (nécessairement non
immergé) à $\mathcal F$ et que $\operatorname{long}(\mathcal F_z)=1$ en son
point générique $z$}. Il est clair que si $\mathcal F$ est un
$\mathcal O_X$-Module cohérent intègre, il est intègre en tout point de $X$~;
inversement, si $\mathcal F$ est \emph{intègre en un point} $x$, il existe un
voisinage ouvert $U$ de $x$ tel que $\mathcal F|U$

